% Generated by D:/tools/scratch_qdd/mmrec/render.py - thread one-liners for the front page.
\begin{itemize}
  \item \textbf{A. Astra 활용} (\ref{rec:A}, 결정 항목 10개): Astra만 할 수 있는 일로 역할을 좁히고(§82), 탐침으로 흐름 상수를 정하고(§86), 한 호출 정확도를 재어 v2 프롬프트를 채택했지만 7--10 cm 어긋남을 못 알아채는 병목(0/13)이 남았다. Astra에 맞춘 단독 조종은 7/7 성공. 유료 호출은 이제 사용자 명시 승인 뒤에만 한다.
  \item \textbf{B. Astra 대체} (\ref{rec:B}, 결정 항목 11개): 목표는 학습 가능한 크기의 오픈 VLM이 Astra 자리를 맡는 것(§96). 영샷은 gpt-5.2 대리 0/4·8B 0/5로 인식에서 막혔고, 무료 시뮬 참값으로 LoRA 학습한 8B는 4/4(접근 오차 125.5 $\rightarrow$ 2.6 mm). 최종 상위는 35B, 8B는 임시. 손으로 준 정보 없이도 되는지(E-STRIP8)와 거리를 모델이 아는 두 방법(E-PT)은 등록 뒤 진행 중.
  \item \textbf{C. 의도 프로토콜과 결합} (\ref{rec:C}, 결정 항목 11개): VLA를 흔한 VLA처럼 쓰지 않고 말단 보정기·세부 조종 스틱으로 두는 틀(§87·§90·§91·§92)을 세우고 결합 배선(Task 1--26)을 끝냈다. 무료 사전 실행에서 VLA 단독이 0/12였고 그 중 8편은 런타임 교착이었다(수정함). 상위$\leftrightarrow$VLA 의도 공유는 필수(§96 보충 2). 남은 것은 학습 --- ser-A-min-3 VLA 재학습, 결합 형식 상위 학습, 결합 폐루프 검증.
  \item \textbf{D. VLA} (\ref{rec:D}, 결정 항목 14개): 움직임 줄만 채택(+0.032)되고 MolmoAct식 보조·카메라 3대·층별 KV·명령 입력은 모두 불채택. 결정 토큰을 청크가 따르지 않는 문제(E-SR0 0.34--0.43)는 시뮬에서 먼 구간 반사실 분기로 0.991까지 풀었지만 실데이터는 0.536(E-SR1e 진행 중). 확신도는 기록만. 다음은 ser-A-min-3 재학습과 끝-끝 학습(E-VLA-E2E).
  \item \textbf{E. 데이터} (\ref{rec:E}, 결정 항목 9개): 조종 표본·복구·자기 인식은 공개 대체가 없어 자체 시뮬(참값 모델)이 주원천이다. 남의 로봇 데이터는 행동 라벨이 아니라 픽셀·QA로 바꿔 쓰는 것이 문헌의 방식이고, 지금 과제는 같은 탁자·같은 물체라 다양한 데이터는 필수(OOD 평가 필수). 공개 데이터 변환은 시제품·관문까지, 좌표계 없는 원천은 픽셀 경로(T0 + T5)가 주 추천. 8B 검증 등록(E-OPEN8·E-XEMB8)은 초안.
  \item \textbf{F. 과정과 규칙} (\ref{rec:F}, 결정 항목 9개): 모든 결정과 실수를 MD·기록책에 남기고(함정 P01--P123), 모든 실험은 사전 등록 + 시작 전·도중·끝 자체 검사, 폐루프 판은 영상 필수, 유료는 명시 승인 뒤. 논문은 가설 기반 완성판 + 빨간 줄 마인드맵 두 판을 4시간마다 갱신. 논문 이름은 PaceNotes(저장소·코드 이름 harvest 유지).
\end{itemize}
